% NOTE (2026-08-30): theory appendix created in the audit-driven restructure.
% Contains the proofs for Section 4 (sec:theory). Style rules: never use the
% em-dash; upright (non-italic) body for theorem-like environments.
\section{Proofs for Section~\ref{sec:theory}}
\label{app:theory}

Throughout this appendix the standing assumptions of the setup paragraph of
Section~\ref{sec:theory} are in force. Because every metric $H\succ0$ makes
$\langle u,Hv\rangle$ an inner product on a finite-dimensional space, all
statements are statements about a Hilbert space, and we work directly in
$\|\cdot\|_H$ without rescaling.

\subsection{Convergence and rate for certified iterations}

\begin{proposition}[Krasnosel'ski\u{\i}--Mann convergence with a
last-iterate rate]
\label{prop:km}
Let $H\succ0$, let $T$ be $\theta$-averaged in $\|\cdot\|_H$ with
$\theta\in(0,1)$, and let $\operatorname{Fix}T\neq\emptyset$. Then the
iteration $z_{k+1}=Tz_k$ converges, from any $z_0$, to some
$z^\star\in\operatorname{Fix}T$; the residual $\|z_{k+1}-z_k\|_H$ is
nonincreasing in $k$; and, for every $k\geq0$,
\begin{equation}
 \|z_{k+1}-z_k\|_H^2
 \;\leq\;
 \frac{\theta}{(1-\theta)\,(k+1)}\;
 \operatorname{dist}_H^2\!\left(z_0,\operatorname{Fix}T\right).
 \label{eq:km-rate}
\end{equation}
\end{proposition}

\begin{proof}
Write $T=(1-\theta)I+\theta S$ with $S$ nonexpansive in $\|\cdot\|_H$; all
norms in this proof are $\|\cdot\|_H$. Fix any
$z^\star\in\operatorname{Fix}T$. Expanding $z_{k+1}=(1-\theta)z_k+\theta
Sz_k$ with the identity
$\|(1-\theta)a+\theta b\|^2=(1-\theta)\|a\|^2+\theta\|b\|^2
-\theta(1-\theta)\|a-b\|^2$, at $a=z_k-z^\star$ and $b=Sz_k-z^\star$, and
bounding $\|Sz_k-z^\star\|\leq\|z_k-z^\star\|$ by nonexpansiveness gives the
Fej\'er inequality
\begin{equation}
 \|z_{k+1}-z^\star\|^2
 \;\leq\;
 \|z_k-z^\star\|^2-\frac{1-\theta}{\theta}\,\|z_{k+1}-z_k\|^2 .
 \label{eq:fejer}
\end{equation}
The map $T$ is nonexpansive, so
$\|z_{k+2}-z_{k+1}\|\leq\|z_{k+1}-z_k\|$. Summing \eqref{eq:fejer} over the
first $k+1$ iterations, keeping only the last (smallest) residual, and
minimizing over $z^\star$ yields \eqref{eq:km-rate}. By \eqref{eq:fejer} the
iterates are bounded and $\|z_{k+1}-z_k\|\to0$, so by continuity every
cluster point of $(z_k)$ is fixed by $T$; reusing \eqref{eq:fejer} with such
a cluster point as $z^\star$ shows $\|z_k-z^\star\|$ is nonincreasing and
vanishes along a subsequence, hence $z_k\to z^\star$. This argument is the
Krasnosel'ski\u{\i}--Mann theorem with its modern rate
\citep[Theorem~1]{RyuYin2022LargeScale}; see also
\citet{Krasnoselskii1955Two}, \citet{Mann1953Mean}, and
\citet{BauschkeCombettes2017Convex}.
\end{proof}

\begin{remark}[Monitored residual, normalization, and summability]
\label{rem:monitored-residual}
The runtime monitor reports $\|z_{k+1}-z_k\|_{H_G}$, the step residual of the
certified map that is actually iterated. For a genome
$\mathrm{relax}_\rho(G')$ the compiled map is
$T_G=(1-\rho)I+\rho T_{G'}$, so
$z_{k+1}-z_k=\rho\,(T_{G'}z_k-z_k)$: the two residuals differ by the factor
$\rho$, which matters when comparing genomes with different relaxation genes,
and \eqref{eq:km-rate} divided by $\rho^2$ bounds the unrelaxed fixed-point
residual $\|T_{G'}z_k-z_k\|_{H_G}^2$. Summing \eqref{eq:fejer} also gives the
summability estimate
$\sum_{k=0}^{\infty}\|z_{k+1}-z_k\|_{H_G}^2\leq
\frac{\theta_G}{1-\theta_G}\operatorname{dist}_{H_G}^2
(z_0,\operatorname{Fix}T_G)$.
\end{remark}

\subsection{Base-scheme certificates}

\paragraph{Reference instantiations.}
The following exact updates fix what the five rows of
Table~\ref{tab:certificates} mean for
$F=f+g+h\circ L$; $\operatorname{prox}_{\alpha q}$ denotes the exact
proximal map of $q$, and $h^*$ is the Fenchel conjugate of $h$. A row is
admissible only if its stated proximal maps can be evaluated. In particular,
writing $g+h\circ L$ as one term does not make its proximal map available.
For FBS, $z=x$ and
$x^+=\operatorname{prox}_{\alpha(g+h\circ L)}(x-\alpha\nabla f(x))$.
For DRS, set $a=g$, $b=f+h\circ L$ and take
$x=\operatorname{prox}_{\alpha a}(z)$,
$y=\operatorname{prox}_{\alpha b}(2x-z)$,
$z^+=z+y-x$; the decoder is $x$.
For Davis--Yin, set $a=g$, $b=h\circ L$ and take
$x=\operatorname{prox}_{\alpha a}(z)$,
$y=\operatorname{prox}_{\alpha b}(2x-z-\alpha\nabla f(x))$,
$z^+=z+y-x$; the decoder is again $x$.
For the primal--dual rows, write $z=(x,u)$ and take
\begin{align*}
 x^+&=\operatorname{prox}_{\tau g}
       \bigl(x-\tau(\nabla f(x)+L^\top u)\bigr),\\
 u^+&=\operatorname{prox}_{\sigma h^*}
       \bigl(u+\sigma L(2x^+-x)\bigr).
\end{align*}
This is Condat--V\~u with decoder $\pi(x,u)=x$; PDHG is the
$\nabla f=0$ specialization. A nonzero $f$ can instead be absorbed into
a proximable primal term, but only when its exact prox is available.
The fixed points of the first three maps decode to minimizers, and those
of the last two are saddle pairs, under the assumptions of
Section~\ref{sec:theory}.

\begin{lemma}[Base schemes]
\label{lem:base-certificates}
For each base scheme of Table~\ref{tab:certificates} with positive stepsizes
satisfying the gate inequality, the pair $(H_G,\theta_G)$ of the table, with
$\theta_G=1/(2-c_G)$, together with the scheme's standard decoder $\pi_G$,
is a construction certificate in the sense of
Definition~\ref{def:certificate}.
\end{lemma}

\begin{proof}
We first establish every row with the true norm $\|L\|$ in place of
$\ell$, then account for the conservative bound.

Since $\nabla f$ is $(1/L_f)$-cocoercive by the Baillon--Haddad theorem
\citep{BaillonHaddad1977Quelques}, the forward step $I-\alpha\nabla f$ is
$(\alpha L_f/2)$-averaged for $\alpha\in(0,2/L_f)$, and the proximal step is
$\tfrac12$-averaged. A composition of $\theta_1$- and $\theta_2$-averaged
maps is $(\theta_1+\theta_2-2\theta_1\theta_2)/(1-\theta_1\theta_2)$-averaged
\citep[Theorem~27]{RyuYin2022LargeScale}, a constant due to
\citet{OguraYamada2002Nonstrictly}; see also
\citet{CombettesYamada2015Compositions}. At $\theta_1=\tfrac12$ and
$\theta_2=c$ it is $1/(2-c)$; this is the FBS row, and the Davis--Yin row is
the same constant \citep[Theorem~28]{RyuYin2022LargeScale}, established by
\citet{DavisYin2017Three}. DRS is
$\tfrac12\bigl(I+R_{\alpha A}R_{\alpha B}\bigr)$, a composition of reflected
resolvents, hence firmly nonexpansive for every $\alpha>0$
\citep[\S2.7]{RyuYin2022LargeScale}, a fact going back to
\citet{LionsMercier1979Splitting}. For the primal--dual rows,
$M_{\tau\sigma}\succ0$ exactly when $\tau\sigma\|L\|^2<1$, and in that metric
PDHG \citep{ChambollePock2011First} is a proximal-point iteration on the
saddle subdifferential, hence firmly nonexpansive in
$\|\cdot\|_{M_{\tau\sigma}}$ \citep[\S3.3]{RyuYin2022LargeScale}; see also
the contraction analysis of \citet{HeYuan2012Convergence}. Condat--V\~u
\citep{Condat2013Primal,Vu2013Splitting} is variable-metric
forward--backward in the same metric \citep{CombettesVu2014Variable}, and its
forward step is $c$-averaged in $\|\cdot\|_{M_{\tau\sigma}}$ with
$c=\tau L_f/\bigl(2(1-\tau\sigma\|L\|^2)\bigr)$, so composition gives
$1/(2-c)$ again \citep[\S3.3]{RyuYin2022LargeScale}. The inequality
$\tau L_f/2+\tau\sigma\|L\|^2<1$ is equivalent, for $\tau,\sigma>0$, to the
conjunction of $\tau\sigma\|L\|^2<1$ and $c<1$, so a single inequality
delivers both (C1) and a positive-definite metric.

Now let $\ell\geq\|L\|$ be the supplied bound and let $c_G$ denote the
tabulated forward-step constant computed from $\ell$. The checked
inequalities in $\ell$ imply the corresponding inequalities in $\|L\|$, so
$H_G\succ0$, and the true forward-step constant $c$ computed from $\|L\|$
satisfies $c\leq c_G<1$, so the compiled iteration is $\theta$-averaged with
$\theta=1/(2-c)$. A $\theta$-averaged map is also $\theta'$-averaged for
every $\theta'\in[\theta,1)$: writing $T=(1-\theta)I+\theta S$ gives
$T=(1-\theta')I+\theta'S'$ with $S'=(1-\theta/\theta')I+(\theta/\theta')S$,
a convex combination of nonexpansive maps. Since
$\theta=1/(2-c)\leq1/(2-c_G)=\theta_G<1$, the pair $(H_G,\theta_G)$
certifies (C1), exactly when $\ell=\|L\|$ and conservatively otherwise.

For (C2) and (C3), no bijection between fixed points and solutions is
needed; it suffices that fixed points exist and that every fixed point
decodes to a solution. For the primal--dual schemes the state is the saddle
pair and, because $M_{\tau\sigma}\succ0$, the fixed points of the
(variable-metric) proximal-point and forward--backward iterations are
exactly the zeros of the saddle operator, i.e.\ the saddle set, which is
nonempty by assumption; $\pi_G$ is the primal coordinate
\citep[\S3.3]{RyuYin2022LargeScale}. For FBS and Davis--Yin the fixed points
are exactly the solutions of the associated inclusion (for Davis--Yin after
applying the decoder $\pi_G=\operatorname{prox}_{\alpha g}$), and for DRS
every fixed point $z$ satisfies
$\operatorname{prox}_{\alpha g}(z)\in\mathcal X^\star$ while every solution
arises from at least one fixed point; these standard correspondences are
\citep[\S2.7]{RyuYin2022LargeScale}. In each case the nonempty saddle set
yields a fixed point through the correspondence, giving (C2), and the
decoder maps all of $\operatorname{Fix}T_G$ into $\mathcal X^\star$, giving
(C3).
\end{proof}

\subsection{Transfer across formulations}

\begin{proposition}[Formulation transfer]
\label{prop:formulation-transfer}
Suppose a transformed composite objective $\widetilde F$ satisfies the
standing assumptions of Section~\ref{sec:theory}, and an exact recovery
map $R$ satisfies
$R(\operatorname*{argmin}\widetilde F)\subseteq
\operatorname*{argmin}F$, where $F$ is the original objective in
\eqref{eq:composite}. If the gate certifies a base scheme and its
relaxation for $\widetilde F$, with decoder $\widetilde\pi_G$, then the
iteration converges in its transformed state space and its limit decodes
via $R\circ\widetilde\pi_G$ to a minimizer of $F$.
\end{proposition}

\begin{proof}
Theorem~\ref{thm:soundness} and Corollary~\ref{cor:convergence} applied
to $\widetilde F$ give convergence to a fixed point whose image under
$\widetilde\pi_G$ minimizes $\widetilde F$. Applying the assumed recovery
map places that minimizer in $\operatorname*{argmin}F$. The gate inequality
and metric must be established anew for the transformed operator; an
equivalent objective alone does not transfer averagedness.
\end{proof}

When relocation moves a constraint into an indicator term, evaluating the
new proximal map requires an exact projection onto that constraint set.
Static row scaling or equilibration likewise requires its own equivalence
and recovery map. The proposition does not identify an approximate or
contract-tested implementation with the exact map.

\subsection{Closure rules}

\begin{lemma}[Relaxation]
\label{lem:relaxation}
Let $T$ be $\theta$-averaged in $\|\cdot\|_H$ and let $\rho>0$ satisfy
$\rho\theta<1$. Then $(1-\rho)I+\rho T$ is $\rho\theta$-averaged in
$\|\cdot\|_H$ with $\operatorname{Fix}\bigl((1-\rho)I+\rho
T\bigr)=\operatorname{Fix}T$.
\end{lemma}

\begin{proof}
Writing $T=(1-\theta)I+\theta S$ with $S$ nonexpansive gives
$(1-\rho)I+\rho T=(1-\rho\theta)I+\rho\theta S$, which is
$\rho\theta$-averaged because $\rho\theta\in(0,1)$, and $\rho>0$ makes the
fixed-point equations equivalent. The decoder and metric are untouched, so
the certificate transports as stated in Theorem~\ref{thm:soundness}(ii).
\end{proof}

The next lemma is a conditional mathematical rule for blends, not an
admission rule of the certified grammar. Its common-fixed-point hypothesis
must be established independently; the current syntactic gate conditions
alone do not establish it.

\begin{lemma}[Blends]
\label{lem:blend}
Let $T_1,T_2$ be $\theta_1$- and $\theta_2$-averaged in a common
$\|\cdot\|_H$ with
$\operatorname{Fix}T_1\cap\operatorname{Fix}T_2\neq\emptyset$, and let
$\lambda\in(0,1)$. Then $T=\lambda T_1+(1-\lambda)T_2$ is
$\bigl(\lambda\theta_1+(1-\lambda)\theta_2\bigr)$-averaged in $\|\cdot\|_H$
and $\operatorname{Fix}T=\operatorname{Fix}T_1\cap\operatorname{Fix}T_2$.
\end{lemma}

\begin{proof}
Write $T_i=(1-\theta_i)I+\theta_iS_i$ and set
$\theta=\lambda\theta_1+(1-\lambda)\theta_2$. Then
$T=(1-\theta)I+\theta\bigl(\tfrac{\lambda\theta_1}{\theta}S_1
+\tfrac{(1-\lambda)\theta_2}{\theta}S_2\bigr)$, and the bracket is a convex
combination of nonexpansive maps, hence nonexpansive; this is the
combination rule of \citet{CombettesYamada2015Compositions}. The inclusion
$\operatorname{Fix}T\supseteq\operatorname{Fix}T_1\cap\operatorname{Fix}T_2$
is immediate. Conversely, let $z\in\operatorname{Fix}T$ and
$w\in\operatorname{Fix}T_1\cap\operatorname{Fix}T_2$. Then
$\|z-w\|\leq\lambda\|T_1z-w\|+(1-\lambda)\|T_2z-w\|\leq\|z-w\|$, so both
inequalities are equalities and $\|T_iz-w\|=\|z-w\|$ for $i=1,2$; the
Fej\'er inequality \eqref{eq:fejer} applied to $T_i$ at $z$ then forces
$\|T_iz-z\|=0$. Hence $z\in\operatorname{Fix}T_1\cap\operatorname{Fix}T_2$.
\end{proof}

\subsection{The Halpern wrapper}

\begin{proposition}[Halpern genomes]
\label{prop:halpern}
Let $T_G$ carry a certificate $(H_G,\theta_G,\pi_G)$ and set
$S_G=I+\theta_G^{-1}(T_G-I)$. Then $S_G$ is nonexpansive in
$\|\cdot\|_{H_G}$ with $\operatorname{Fix}S_G=\operatorname{Fix}T_G$, and the
compiled iteration
\begin{equation}
 z_{k+1}=\lambda_{k+1}z_{\rm anc}+(1-\lambda_{k+1})S_Gz_k ,
 \label{eq:halpern}
\end{equation}
under any admitted schedule, i.e.\ $\lambda_k\to0$,
$\sum_k\lambda_k=\infty$, and $\sum_k|\lambda_{k+1}-\lambda_k|<\infty$,
converges to $P^{H_G}_{\operatorname{Fix}T_G}(z_{\rm anc})$, the projection
of the anchor onto $\operatorname{Fix}T_G$ in $\|\cdot\|_{H_G}$
\citep{Halpern1967Fixed,Wittmann1992Approximation}. The canonical schedule
$\lambda_{k+1}=1/(k+2)$ meets all three conditions and, when the wrapper
anchors at the initial iterate ($z_{\rm anc}=z_0$), in addition gives
\begin{equation}
 \|S_Gz_k-z_k\|_{H_G}^2
 \;\leq\;
 \frac{4\operatorname{dist}_{H_G}^2\!\left(z_0,\operatorname{Fix}
 T_G\right)}{(k+1)^2}
 \label{eq:halpern-rate}
\end{equation}
\citep{Lieder2021Optimized,Kim2021Accelerated}.
\end{proposition}

\begin{proof}
$S_G$ is precisely the nonexpansive map in the averaged decomposition
$T_G=(1-\theta_G)I+\theta_GS_G$ of (C1), and $T_Gz=z$ if and only if
$S_Gz=z$ because $\theta_G>0$. Since $(\mathcal Z_G,\langle\cdot,H_G
\cdot\rangle)$ is a Hilbert space, the cited Halpern results apply verbatim
in $\|\cdot\|_{H_G}$: convergence to the projection of the anchor is
\citet{Wittmann1992Approximation}, and the rate \eqref{eq:halpern-rate} is
\citet{Lieder2021Optimized}.
\end{proof}

Rate \eqref{eq:halpern-rate} is an $O(1/k^2)$ bound on the same squared
residual that \eqref{eq:km-rate} bounds at $O(1/k)$, and it is exactly
optimal in the following sense: \citet{ParkRyu2022Exact} exhibit a matching
lower bound for deterministic fixed-point methods that access $S_G$ only
through evaluations, with iterates constrained to the span of the anchor and
past evaluations. The claim is optimality within that oracle class, not an
unrestricted impossibility statement.

\subsection{Safeguarded schedules}

\begin{lemma}[Summable-error safeguard]
\label{lem:safeguard}
Let $T$ be $\theta$-averaged in $\|\cdot\|_H$ with
$\operatorname{Fix}T\neq\emptyset$, and let $z_{k+1}=Tz_k+e_k$ with
$\|e_k\|_H\leq\eta_k$ and $\sum_k\eta_k<\infty$. Then $(z_k)$ is
quasi-Fej\'er monotone with respect to $\operatorname{Fix}T$ in
$\|\cdot\|_H$ and converges to a point of $\operatorname{Fix}T$.
\end{lemma}

\begin{proof}
For any $z^\star\in\operatorname{Fix}T$, nonexpansiveness gives the one-step
inequality
$\|z_{k+1}-z^\star\|_H\leq\|Tz_k-z^\star\|_H+\|e_k\|_H
\leq\|z_k-z^\star\|_H+\eta_k$ with $\sum_k\eta_k<\infty$, which is
quasi-Fej\'er monotonicity; in particular
$\|z_k-z^\star\|_H\leq M:=\|z_0-z^\star\|_H+\sum_j\eta_j$ for every $k$.
Combining the one-step inequality with the Fej\'er
inequality~\eqref{eq:fejer} for the exact step gives
\begin{equation*}
 \|z_{k+1}-z^\star\|_H^2
 \;\leq\;
 \|z_k-z^\star\|_H^2-\frac{1-\theta}{\theta}\,\|Tz_k-z_k\|_H^2
 +2M\eta_k+\eta_k^2 ,
\end{equation*}
and summing shows $\sum_k\|Tz_k-z_k\|_H^2<\infty$, hence
$\|Tz_k-z_k\|_H\to0$. The iterates are bounded, so by continuity some
cluster point $z^\star$ is fixed by $T$; quasi-Fej\'er monotonicity makes
$\|z_k-z^\star\|_H+\sum_{j\geq k}\eta_j$ nonincreasing, so
$\|z_k-z^\star\|_H$ converges, and it vanishes along a subsequence, hence
$z_k\to z^\star$. This is the classical inexact
Krasnosel'ski\u{\i}--Mann analysis
\citep{Combettes2001QuasiFejerian,Combettes2004Solving}.
\end{proof}

The same mechanism underlies the inexact analysis of Douglas--Rachford and
ADMM \citep{EcksteinBertsekas1992Douglas} and safeguarded Anderson
acceleration \citep{Zhang2020Globally}, the design template for the
$\mathrm{switch}$ gene. The guarantee is convergence, not a rate: the budget
buys back the limit an aggressive step would otherwise forfeit.

\FloatBarrier
